% TEMPLATE-MILC-v3.tex - plantilla maestra UNICA de las guias MILC v3.
%
% La usan las cuatro familias de guias, cada una con su driver:
%   · Grados 6-11          -> scripts/build-guias-g11.py
%   · Bebras               -> scripts/build-guias-bebras.py
%   · Semillero            -> scripts/build-guias-semillero.py
%   · Territorio interior  -> scripts/build-guias-territorio-interior.py
%
% Antes habia tres copias casi identicas (~400 lineas iguales, 6 distintas):
% cada mejora habia que aplicarla tres veces, y en la practica no se hacia
% (el motor de recursos vivia solo en dos de ellas). Lo que varia entre
% familias esta parametrizado en 6 placeholders que cada driver llena
% (se nombran aqui SIN su sintaxis real: el builder reemplaza por texto plano,
% tambien dentro de los comentarios, y un valor multilinea rompe el preambulo):
%
%   HEADER_IZQ        encabezado izquierdo de pagina
%   HEADER_DER        encabezado derecho de pagina
%   PORTADA_ETIQUETA  rotulo sobre el numero grande (GRADO/MOMENTO/MODULO)
%   PORTADA_SERIE     linea de serie de la portada
%   PORTADA_ID        identificador de la guia en la portada
%   PORTADA_CIRCULO   el circulito de grado (°), o vacio si no aplica
%   PDF_TITULO        titulo en texto plano (sin \textbf ni comillas TeX) para
%                     los metadatos del PDF (/Title); lo produce lib_xelatex.texto_pdf
%
% Composicion (auditoria 2026-09-03): las bandas de estacion piden espacio
% (\Needspace*) y prohiben el salto justo despues (\nopagebreak), asi no quedan
% huerfanas al pie; las cajas largas (softbox, drawbox, infoband, puentebox) son
% `breakable`, asi una caja de media pagina ya no arrastra todo a la siguiente
% dejando la anterior al 40 %. Todo texto sobre mostaza o verde va en milcNegro
% (blanco sobre #E5B400 da 1,93:1 y sobre #58A923 2,95:1; ninguno pasa WCAG AA).
% El resto de claves las documenta content/guias/_SCHEMA.md. El script valida
% que no queden placeholders sin reemplazar antes de escribir el archivo final.

% Etiquetado PDF (latex-lab): /Lang, estructura y texto alternativo de las
% figuras, para lectores de pantalla. Debe ir ANTES de \documentclass.
\DocumentMetadata{lang=es-CO,tagging=on}
\documentclass[11pt,letterpaper]{article}
% headheight=24pt: la cabecera derecha lleva dos lineas (identidad de la guia
% y estacion actual). headsep se acorta en la misma medida para que el bloque
% de texto quede exactamente donde estaba con headheight=12pt.
\usepackage[margin=2.0cm,headheight=24pt,headsep=13pt]{geometry}
\usepackage[table]{xcolor}
\usepackage{fontspec}
\usepackage[spanish]{babel}
\usepackage{tikz}
% Librerías TikZ para los diagramas de las guías (bloque `recursos:`):
%  · babel  → babel en español vuelve activos caracteres como «>», y TikZ los usa
%             en sus opciones (>=stealth, ->). Sin esta librería, cualquier
%             diagrama con flechas revienta con «\language@active@arg> has an
%             extra }». Debe ir ANTES de que se dibuje nada.
%  · positioning        → sintaxis «below left=2pt and 3pt».
%  · decorations.pathreplacing → llaves ({brace}) para señalar rangos.
\usetikzlibrary{babel,positioning,decorations.pathreplacing}
\usepackage{graphicx}
% Circuitos: el trabajo de electrónica se hace hoy con micro:bit y sus sensores
% integrados (MakeCode, sin cableado), así que no se necesitan esquemáticos.
% Si algún día entran montajes con protoboard, basta descomentar esta línea:
% los diagramas .tex/.tikz declarados en `recursos:` ya se incrustan solos.
% \usepackage{circuitikz}
\usepackage[most]{tcolorbox}
\usepackage{tabularx}
\usepackage{array}
\usepackage{fancyhdr}
\usepackage{titlesec}
\usepackage[document]{ragged2e}  % bandera a la derecha en todo el cuerpo
\usepackage{enumitem}
\usepackage{microtype}
\usepackage{etoolbox}
\usepackage{needspace}
% hyperref al final del preambulo: metadatos del PDF (titulo, autor, idioma,
% familia), marcadores por estacion (\pdfbookmark) y enlaces sin recuadro.
\usepackage[hidelinks,
  pdftitle={Fórmulas básicas — SUMA, PROMEDIO, MAX, MIN, CONTAR},
  pdfauthor={PhD. Álvaro Cárdenas Orozco},
  pdfsubject={Tecnología e Informática · Grado 9 · Guía 23 · MILC v3},
  pdflang={es-CO},
  bookmarksopen=true,bookmarksnumbered=false]{hyperref}

% Helvetica Neue no trae la flecha U+2192, asi que cada «→» escrito literalmente
% en el YAML se compilaba a NADA, en silencio (462 flechas en 61 guias: rutas del
% tipo «paleta → tipografia → lockup» salian rotas). Mapeamos el caracter a su
% version matematica, que si tiene glifo. Asi el autor puede seguir escribiendo
% «→» comodamente en el YAML.
\usepackage{newunicodechar}
\newunicodechar{→}{\ensuremath{\rightarrow}}
% Mismo problema con los simbolos de marca y vinetas: el «✓» de «una palomita ✓»
% salia como cuadrito vacio en el PDF de todas las guias que lo usaban.
\usepackage{pifont}
\newunicodechar{✓}{\ding{51}}
\newunicodechar{✗}{\ding{55}}
\newunicodechar{✔}{\ding{52}}
\newunicodechar{•}{\textbullet}
\newunicodechar{≥}{\ensuremath{\geq}}
\newunicodechar{≤}{\ensuremath{\leq}}

\setmainfont{Helvetica Neue}
\setsansfont{Helvetica Neue}
\setlength{\parindent}{0pt}
\setlength{\parskip}{6pt}
% Sin particion de palabras: en 6.º-7.º una palabra cortada por guion al final
% de linea («Comuni-cacion») frena la lectura mucho mas que una linea corta.
\hyphenpenalty=10000
\exhyphenpenalty=10000
\pretolerance=2000
\tolerance=3000
\setlength{\emergencystretch}{3em}
\linespread{1.10}
\renewcommand{\arraystretch}{1.25}
\setlist{leftmargin=5mm,itemsep=1mm,topsep=1mm}
\newcolumntype{Y}{>{\RaggedRight\arraybackslash}X}

\definecolor{milcMagenta}{HTML}{E6007E}
\definecolor{milcVino}{HTML}{5A0038}
\definecolor{milcTurquesa}{HTML}{0093A5}
\definecolor{milcVerde}{HTML}{58A923}
\definecolor{milcMostaza}{HTML}{E5B400}
\definecolor{milcOcre}{HTML}{B86600}
\definecolor{milcNegro}{HTML}{241816}
\definecolor{milcGris}{HTML}{F7F7F4}
\definecolor{milcLinea}{HTML}{E8E2DD}
\definecolor{milcGradePrimary}{HTML}{7C3AED}
\definecolor{milcGradeDark}{HTML}{4C1D95}
\definecolor{milcGradeSoft}{HTML}{EDE9FE}
\definecolor{milcGradeAccent}{HTML}{CFE7FF}
\definecolor{milcGradeText}{HTML}{FFFFFF}
\color{milcNegro}

\pagestyle{fancy}
\fancyhf{}
% Cabecera de dos lineas: identidad de la guia arriba y, a la derecha y en
% gris, la estacion en curso (\markright desde \milcestacion). Antes de la
% primera estacion la segunda linea va vacia; el \strut mantiene la altura
% para que la linea de arriba no baile entre paginas.
\lhead{\small Tecnología e Informática · Grado 9\\[-1pt]{\footnotesize\strut}}
\rhead{\small Guía 23 · MILC v3\\[-1pt]{\footnotesize\strut\color{milcNegro!65}\rightmark}}
\cfoot{\small\thepage}
\renewcommand{\headrulewidth}{0pt}

% Sobre mostaza (#E5B400) y verde (#58A923) el texto va en milcNegro; sobre el
% resto de la paleta, en blanco. Se usa dentro de fonttitle/fontupper, donde un
% \color posterior manda sobre coltitle.
\newcommand{\milctintasobre}[1]{%
  \ifboolexpr{test{\ifstrequal{#1}{milcMostaza}} or test{\ifstrequal{#1}{milcVerde}}}%
    {\color{milcNegro}}{\color{white}}}

\newtcolorbox{titlebox}[1]{
  enhanced,colback=#1,colframe=#1,arc=7mm,boxrule=0pt,
  left=7mm,right=7mm,top=3mm,bottom=3mm,
  before skip=8pt,after skip=8pt,
  fontupper=\sffamily\bfseries\Large\color{white}
}

\newtcolorbox{softbox}[3]{
  enhanced,breakable,pad at break=3mm,
  colback=#2,colframe=#1,borderline west={4pt}{0pt}{#1},
  arc=2mm,boxrule=.4pt,left=4mm,right=4mm,top=3mm,bottom=3mm,
  before skip=6pt,after skip=8pt,title={#3},coltitle=white,
  colbacktitle=#1,fonttitle=\sffamily\bfseries\milctintasobre{#1},fontupper=\RaggedRight,
  attach boxed title to top left={xshift=3mm,yshift=-2mm},
  boxed title style={arc=2mm, boxrule=0pt}
}

\newtcolorbox{drawbox}[2]{
  enhanced,breakable,pad at break=4mm,
  colback=white,colframe=#1,arc=4mm,boxrule=1pt,
  left=5mm,right=5mm,top=4mm,bottom=4mm,before skip=8pt,after skip=8pt,
  title={#2},coltitle=white,colbacktitle=#1,fonttitle=\sffamily\bfseries\milctintasobre{#1},
  fontupper=\RaggedRight,
  attach boxed title to top left={xshift=4mm,yshift=-2mm},
  boxed title style={arc=3mm, boxrule=0pt}
}

\newtcolorbox{infoband}[1]{
  enhanced,breakable,pad at break=2mm,colback=#1!8,colframe=#1!50,arc=2mm,boxrule=.5pt,
  left=4mm,right=4mm,top=2mm,bottom=2mm,before skip=4pt,after skip=4pt,
  fontupper=\small\RaggedRight
}

% Puente narrativo entre secciones (contrato editorial Conéctate).
% Da continuidad: "Acabas de X. Ahora vas a Y porque Z."
% Visualmente: banda izquierda en color del grado, texto en itálica pequeña.
\newtcolorbox{puentebox}{
  enhanced,breakable,pad at break=2mm,colback=milcGradeSoft,colframe=milcGradeDark,
  borderline west={3pt}{0pt}{milcGradeDark},
  arc=0mm,boxrule=0pt,left=5mm,right=4mm,top=2.5mm,bottom=2.5mm,
  before skip=4pt,after skip=8pt,
  fontupper=\itshape\small\color{milcNegro}\RaggedRight
}
% La flecha va en modo matematico a proposito: Helvetica Neue NO tiene el glifo
% U+2192, asi que \textrightarrow se compilaba a NADA («Missing character: There
% is no → in font Helvetica Neue») y el puente salia sin flecha. En modo
% matematico la flecha viene de la fuente matematica y si se dibuja.
\newcommand{\puente}[1]{%
  \begin{puentebox}%
    \textcolor{milcGradeDark}{$\rightarrow$}\hspace{2mm}#1%
  \end{puentebox}%
}

\newcommand{\linead}[1]{\noindent\color{milcLinea}\rule{#1}{0.5pt}\color{milcNegro}}
\newcommand{\respuesta}[1]{\par\vspace{2mm}\foreach \n in {1,...,#1}{\noindent\color{milcLinea}\rule{\linewidth}{0.5pt}\par\vspace{5mm}}\color{milcNegro}}
\newcommand{\checkbox}{\textcolor{milcGradeDark}{\fbox{\phantom{x}}}\hspace{5pt}}

% ─── Listas aireadas para las guías socioemocionales ────────────────────
% Componer las verificaciones como «\checkbox texto\\» deja el texto que salta
% de línea debajo del cuadrito y apelmaza el bloque. Estos entornos alinean la
% continuación y airean los ítems. Los usa Territorio Interior; cualquier
% familia puede hacerlo.
\newlist{checklista}{itemize}{1}
\setlist[checklista]{
  label=\textcolor{milcGradeDark}{\fbox{\phantom{x}}},
  leftmargin=9mm, labelsep=3.2mm, itemsep=2.4mm,
  topsep=2.5mm, parsep=0pt, partopsep=0pt, before=\RaggedRight
}
\newlist{pasoslista}{enumerate}{1}
\setlist[pasoslista]{
  label=\textcolor{milcGradeDark}{\sffamily\bfseries\arabic*.},
  leftmargin=8mm, labelsep=3mm, itemsep=2.8mm,
  topsep=1mm, parsep=0pt, partopsep=0pt, before=\RaggedRight
}

\newcommand{\phasecard}[4]{
\begin{tcolorbox}[enhanced,colback=#3,colframe=#2,arc=3mm,boxrule=.5pt,height=3.05cm,valign=center,halign=center,left=1.5mm,right=1.5mm,top=2mm,bottom=2mm]
{\sffamily\bfseries\fontsize{22}{24}\selectfont\textcolor{#2}{#1}}\par
{\sffamily\bfseries\small #4}
\end{tcolorbox}
}

% ── Piezas del rediseno 2026 (legibilidad 6.º-11.º) ───────────────────
% Banda de estacion: reemplaza el titlebox macizo. Titulo a la izquierda,
% dato practico (tiempo/modalidad) a la derecha, en una sola linea.
% Nunca huerfana: pide 9 lineas libres (o salta de pagina rellenando con
% \vfill) y prohibe el salto justo despues de la banda. Nueve y no seis:
% la caja partible que sigue necesita ~80pt (titulo adosado, relleno y dos
% lineas) para arrancar en la misma pagina; con menos, tcolorbox la manda
% entera a la siguiente y la banda se queda sola. El penalty 10000 va
% en `after`, ANTES del espacio posterior: un `after skip` (aunque sea 0pt)
% es un \vskip tras la caja, y ese glue es un punto de corte legal que un
% \nopagebreak escrito despues de \end{tcolorbox} ya no cierra. Cada banda
% deja marcador PDF y actualiza la cabecera derecha.
\newcounter{milcestacion}
\newcommand{\milcestacion}[2]{%
\Needspace*{9\baselineskip}%
\stepcounter{milcestacion}%
\pdfbookmark[1]{#2}{milcestacion\themilcestacion}%
\markright{#2}%
\begin{tcolorbox}[enhanced,colback=#1,colframe=#1,arc=2mm,boxrule=0pt,
  left=5mm,right=5mm,top=2.6mm,bottom=2.6mm,before skip=11pt,
  after={\par\nopagebreak\vskip7pt}]
{\sffamily\bfseries\fontsize{15}{18}\selectfont\milctintasobre{#1}#2}
\end{tcolorbox}}

% Tarjeta de paso numerada: sustituye las tablas «Pilar / Aplicacion», que
% eran formato de consulta docente, no de estudiante. El numero va en un
% circulo sobre el borde para que la secuencia se lea de un vistazo.
\newcommand{\milcpaso}[3]{%
\Needspace{4\baselineskip}%
\begin{tcolorbox}[enhanced,colback=white,colframe=milcLinea,arc=1.5mm,boxrule=.9pt,
  left=14mm,right=4mm,top=3mm,bottom=3mm,before skip=4pt,after skip=4pt,
  fontupper=\RaggedRight,
  overlay={\node[anchor=north west,circle,fill=#2,text=white,inner sep=0pt,
    minimum size=7.4mm,font=\sffamily\bfseries\small]
    at ([xshift=3.6mm,yshift=-3.2mm]frame.north west) {#1};}]
#3
\end{tcolorbox}}

% Casilla marcable de tamano util con lapiz (la anterior era \fbox{\phantom{x}},
% de alto variable segun el texto que la seguia).
\newcommand{\milcmarca}{\framebox[3.8mm]{\rule{0pt}{3.4mm}}}

% Fila marcable de lista suelta (checks de la estacion 1).
\newcommand{\milccheckrow}[1]{%
\par\vspace{1.8mm}\noindent
\makebox[6.4mm][l]{\raisebox{-.55mm}{\textcolor{milcNegro}{\milcmarca}}}%
\parbox[t]{\dimexpr\linewidth-6.4mm\relax}{\RaggedRight #1}\par}

% Celda marcable dentro de la rubrica (tabla de autoevaluacion). Misma
% sangria francesa, pero pensada para vivir dentro de una columna X.
\newcommand{\milcopcion}[1]{%
\makebox[5.2mm][l]{\raisebox{-.55mm}{\textcolor{milcNegro}{\milcmarca}}}%
\parbox[t]{\dimexpr\linewidth-5.2mm\relax}{\RaggedRight #1}}

% ── Recursos visuales de la guía ──────────────────────────────────────
% Declarados en el bloque `recursos:` del YAML y emitidos por el builder.
% \guiaFigura   → imagen rasterizada o PDF (foto, carta celeste, montaje).
% \guiaDiagrama → fuente TikZ/circuitikz (.tex) incrustada como vector:
%                 es la vía para circuitos y esquemáticos editables.
% [#1] = texto alternativo (alt) · #2 = ruta relativa al .tex · #3 = pie
% El alt llega al PDF etiquetado (/Alt) y lo lee el lector de pantalla.
\newcommand{\guiaFigura}[3][]{%
\begin{center}
\includegraphics[width=0.88\linewidth,height=0.38\textheight,keepaspectratio,alt={#1}]{#2}\\[1.5mm]
{\footnotesize\itshape\textcolor{milcNegro!75}{#3}}
\end{center}
}

\newcommand{\guiaDiagrama}[2]{%
\begin{center}
\input{#1}\\[1.5mm]
{\footnotesize\itshape\textcolor{milcNegro!75}{#2}}
\end{center}
}

\begin{document}
\thispagestyle{empty}

% ── Portada ───────────────────────────────────────────────────────────
% Antes: un rectangulo de color a pagina completa. Se veia bien en pantalla,
% pero en la fotocopiadora del colegio se comia el toner y salia gris sucio.
% Ahora la banda ocupa el tercio superior y el resto va en blanco.
\begin{tikzpicture}[remember picture,overlay]
  \path[fill=milcGradePrimary]
    (current page.north west) rectangle ([yshift=-10.6cm]current page.north east);

  \node[anchor=north west,text=milcGradeText] at ([xshift=2.0cm,yshift=-1.55cm]current page.north west)
    {\sffamily\bfseries\fontsize{12}{14}\selectfont GRADO};
  \node[anchor=north west,text=milcGradeText,opacity=.85] at ([xshift=2.0cm,yshift=-2.35cm]current page.north west)
    {\sffamily\bfseries\fontsize{8.5}{10}\selectfont SERIE GUÍAS · TECNOLOGÍA E INFORMÁTICA};
  \node[anchor=north east,text=milcGradeText,opacity=.85,align=right] at ([xshift=-2.0cm,yshift=-1.55cm]current page.north east)
    {\sffamily\bfseries\fontsize{9}{12}\selectfont I.E. Sor María Juliana\\Cartago, Valle del Cauca};

  \node[anchor=west,text=milcGradeText] (gradonum) at ([xshift=1.85cm,yshift=-7.1cm]current page.north west)
    {\sffamily\bfseries\fontsize{112}{112}\selectfont 9};
  % El circulito de grado (°) solo aplica a las guias de grado («11°»); en
  % Bebras y Semillero el numero grande es un momento/modulo, no un grado.
  % Se posiciona relativo al nodo `gradonum`, no en coordenadas absolutas.
  \draw[milcGradeText,line width=1.6pt]
    ([xshift=2.0mm,yshift=-4.5mm]gradonum.north east) circle (.15cm);

  \node[anchor=west,text=milcGradeText,text width=11.4cm,align=left] at ([xshift=1.5cm,yshift=0.5cm]gradonum.east)
    {
     {\sffamily\bfseries\fontsize{9}{11}\selectfont\MakeUppercase{Período 3 · Datos --- del registro al insight}}\\[3mm]
     {\sffamily\bfseries\fontsize{21}{25}\selectfont\setlength{\baselineskip}{27pt}Fórmulas básicas ---\\[4pt]cinco preguntas\\[4pt]a la tabla}\\[3.5mm]
     {\sffamily\bfseries\fontsize{15}{18}\selectfont Guía 23-9-TIC}
    };
\end{tikzpicture}

\vspace*{9.4cm}
\vfill

{\sffamily\bfseries\fontsize{8}{10}\selectfont\textcolor{milcGradeDark}{LO QUE TE LLEVAS HOY}}

\begin{tcolorbox}[enhanced,colback=white,colframe=milcNegro,arc=2mm,boxrule=1.6pt,
  left=5mm,right=5mm,top=4mm,bottom=4mm,before skip=3pt,after skip=12pt,
  fontupper=\RaggedRight\fontsize{11}{15.5}\selectfont]
Cinco indicadores calculados con fórmula sobre tu propia tabla: total, promedio, máximo, mínimo y cantidad de casos. Cada uno con su sintaxis escrita, su resultado verificado y una frase que diga qué significa en tu contexto, no en general. Cierra señalando cuál de los cinco es el que de verdad sirve para la decisión que tenías en mente.
\end{tcolorbox}

\begin{tcolorbox}[enhanced,colback=milcGris,colframe=milcGris,arc=2mm,boxrule=0pt,
  left=5mm,right=5mm,top=3.5mm,bottom=3.5mm,after skip=12pt]
{\sffamily\bfseries\fontsize{8}{10}\selectfont\textcolor{milcNegro!55}{ESTA GUÍA ES DE}}\\[3mm]
\begin{tabularx}{\linewidth}{X p{3.2cm} p{3.2cm}}
\linead{\linewidth} & \linead{3.0cm} & \linead{3.0cm}\\[-1mm]
{\fontsize{8.5}{10}\selectfont\textcolor{milcNegro!55}{Nombre completo}} &
{\fontsize{8.5}{10}\selectfont\textcolor{milcNegro!55}{Grupo}} &
{\fontsize{8.5}{10}\selectfont\textcolor{milcNegro!55}{Fecha}}\\
\end{tabularx}
\end{tcolorbox}

\vfill

\noindent\textcolor{milcLinea}{\rule{\linewidth}{1pt}}\\[2mm]
{\sffamily\bfseries\fontsize{9.5}{12}\selectfont Autor: Dr. Álvaro Cárdenas Orozco}\\[1mm]
{\sffamily\fontsize{8.5}{11}\selectfont\textcolor{milcNegro!55}{Metodología MILC · apertura ancestral · pensamiento computacional · triángulo Dussel–estoicismo–Floridi}}

\newpage

\section*{Apertura: Fórmulas básicas --- SUMA, PROMEDIO, MAX, MIN, CONTAR}

\begin{softbox}{milcGradeDark}{milcGradeSoft}{Saber ancestral}
En una tienda de barrio del Valle del Cauca hay un cuaderno. La vendedora anota ahí quién debe y cuánto, y va tachando a medida que le pagan. Nadie firma nada: al fiado se entra demostrando que uno es «buena paga» y se sale cuando uno deja de serlo (Martínez Benavides, 2021). Fíjate en lo que hace ese cuaderno. La palabra vale porque hay memoria, y el registro vale porque cualquiera puede ponerlo a prueba: se abre, se suma y se compara. Es una hoja de cálculo hecha a mano, con las mismas operaciones que vas a usar hoy. La \textbf{cara de exclusión} es doble. La sanción por incumplir no es proporcional ni apelable: es la vergüenza en público, delante de testigos. Y a quien la tendera no conoce, no le fía, así que el circuito solo cuenta a los que ya estaban adentro.\par\smallskip{\footnotesize\textcolor{milcNegro!70}{\textbf{Fuente.} Martínez Benavides, A. (2021). Circuitos crediticios: fiado y trabajo relacional en un pequeño negocio en Cali, Colombia. \emph{Estudios Sociológicos de El Colegio de México, 39}(116), 467--494. https://doi.org/10.24201/es.2021v39n116.1946}}
\end{softbox}

\begin{softbox}{milcTurquesa}{milcGris}{Saber contemporáneo}
Cinco fórmulas responden casi todo lo que una tabla puede contestar por sí sola. \texttt{SUMA} responde cuánto hay en total. \texttt{PROMEDIO} responde cuánto suele haber por caso. \texttt{MAX} y \texttt{MIN} responden dónde están los extremos. \texttt{CONTAR} responde cuántos casos hay. Las cinco piden lo mismo para funcionar: una columna con un solo tipo de dato y un rango bien seleccionado. Y ahí está el error más común: seleccionar celda por celda en vez de tomar el rango completo, porque si mañana llega una fila nueva, la fórmula no la ve. Pero la parte que decide no es la sintaxis. Es la pregunta previa: qué vas a decidir con ese número. Un promedio de notas sirve para saber si hay que reforzar; un máximo sirve para saber si alguien se quedó sin dormir. No son intercambiables.
\end{softbox}

\begin{softbox}{milcMostaza}{milcGris}{Pregunta-puente}
\textit{El cuaderno de la tienda vale porque cualquiera puede abrirlo, sumar y comparar. Tu hoja de cálculo, ¿aguantaría que alguien la revisara así?}
\end{softbox}

\begin{softbox}{milcVerde}{milcGris}{Saber y saber hacer}
Al terminar podrás: (1) \textbf{identificar} cuánto se desvían tus cálculos mentales de los de la hoja y dónde estuvo la diferencia; (2) \textbf{explicar} qué pregunta responde cada una de las cinco fórmulas y qué exige para funcionar; (3) \textbf{aplicar} las cinco sobre tu tabla e interpretar cada resultado en tu contexto.
\end{softbox}



\small
\begin{tabularx}{\textwidth}{>{\bfseries}p{3.0cm}Y}
\rowcolor{milcGradePrimary!12}Autor & Dr. Álvaro Cárdenas Orozco\\
DBA & Tratamiento y visualización honesta de datos para la toma de decisiones (MEN, T\&I)\\
Referentes & Ética del dato · Visualización honesta · Pensamiento computacional\\
Producto final & Cinco indicadores calculados con fórmula sobre tu propia tabla: total, promedio, máximo, mínimo y cantidad de casos. Cada uno con su sintaxis escrita, su resultado verificado y una frase que diga qué significa en tu contexto, no en general. Cierra señalando cuál de los cinco es el que de verdad sirve para la decisión que tenías en mente.\\
\end{tabularx}
\normalsize

\puente{En la sesión 1 recogiste veinte filas y en la 2 las pusiste en forma. Hoy le haces a esa tabla las cinco preguntas básicas. En la sesión 4 aprendes a preguntarle solo por una parte, con filtros.}

\milcestacion{milcGradeDark}{Ruta MILC: así vamos a aprender}
\begin{tabularx}{\textwidth}{>{\centering\arraybackslash}X>{\centering\arraybackslash}X>{\centering\arraybackslash}X>{\centering\arraybackslash}X}
\phasecard{1}{milcVerde}{milcGris}{Escucha\\[-1mm]\small Observo, escucho y pregunto.} &
\phasecard{2}{milcTurquesa}{milcGris}{Entender\\[-1mm]\small Sistematizo y ordeno.} &
\phasecard{3}{milcMagenta}{milcGris}{Hacer\\[-1mm]\small Construyo y aplico.} &
\phasecard{4}{milcVino}{milcGris}{Cierre\\[-1mm]\small Evalúo y reflexiono.}
\end{tabularx}


\puente{Antes de la teoría vas a calcular de cabeza y después con fórmula. La diferencia entre los dos resultados es la clase de hoy.}

\milcestacion{milcVerde}{Estación 1 de 3 · Escucha}

\begin{softbox}{milcVerde}{milcGris}{Escena para escuchar}
\textbf{Actividad 1 · IDENTIFICA --- De cabeza y con fórmula} (15 min · individual). (1) Abre tu tabla de la sesión 2, con su columna numérica. (2) Calcula mentalmente el total, el promedio, el máximo, el mínimo y cuántas filas hay. (3) Anota tus cinco respuestas y cuánto tardaste. (4) Calcula ahora los cinco con fórmula y anota los resultados. (5) Compara y marca en cuál te alejaste más, y escribe por qué crees que fue ese.
\end{softbox}

{\sffamily\bfseries\fontsize{8}{10}\selectfont\textcolor{milcNegro!55}{MARCA CUANDO LO HAYAS HECHO}}
\milccheckrow{Tengo los cinco resultados de cabeza y los cinco con fórmula.}
\milccheckrow{Marqué en cuál me alejé más.}
\milccheckrow{Escribí por qué creo que fue ese y no otro.}

\begin{infoband}{milcVerde}


\textbf{Lo que va al cuaderno (Actividad 1 · IDENTIFICA).} \textbf{Título:} «Actividad 1 --- De cabeza y con fórmula». \textbf{Formato:} tabla de 5 filas y 3 columnas (indicador / de cabeza / con fórmula), con marca en la mayor diferencia. \textbf{Extensión:} un tercio de página. \textbf{Sabes que terminaste cuando} tienes escrito por qué falló ese cálculo y no otro.
\end{infoband}

\puente{Ya viste dónde falla tu cálculo mental. Ahora vas a saber qué pregunta responde cada fórmula y qué necesita para no mentir.}

\milcestacion{milcTurquesa}{Estación 2 de 3 · Entender}

\begin{softbox}{milcTurquesa}{milcGris}{Lo que vas a entender}
\textbf{Actividad 2 · EXPLICA --- Qué pregunta responde cada una} (20 min · parejas). (1) Con tu pareja, escriban las cinco fórmulas con la pregunta que responde cada una, en lenguaje corriente. (2) Escriban qué exige cada una para funcionar. (3) Tomen un dato de sus tablas y digan qué indicador lo describe mejor y cuál lo describe peor. (4) Busquen un caso donde el promedio esconda algo y escriban qué indicador lo mostraría.
\end{softbox}

\milcpaso{1}{milcTurquesa}{\textbf{Las cinco preguntas.} SUMA: cuánto hay en total. PROMEDIO: cuánto suele haber por caso. MAX y MIN: dónde están los extremos. CONTAR: cuántos casos hay. Cinco preguntas distintas y ningún atajo entre ellas.}
\milcpaso{2}{milcTurquesa}{\textbf{Rango antes que celda suelta.} Selecciona la columna completa, no celda por celda. Si mañana llega una fila nueva, la fórmula con rango la incluye sola y la otra no, y nadie te avisa del error.}
\milcpaso{3}{milcTurquesa}{\textbf{El promedio esconde a los extremos.} Un promedio de cuatro horas de pantalla puede tapar a alguien con once. Por eso el máximo y el mínimo no son adornos: son la parte del reporte que muestra a quién le está pasando algo.}
\milcpaso{4}{milcTurquesa}{\textbf{La pregunta va antes que la fórmula.} Qué vas a decidir con ese número. Un promedio sirve para saber si hay que reforzar; un máximo, para saber si alguien se quedó sin dormir. No son intercambiables.}

\begin{drawbox}{milcTurquesa}{Las cinco fórmulas}
\texttt{SUMA} \textbf{→} cuánto en total. \\[1mm] \texttt{PROMEDIO} \textbf{→} cuánto por caso. \\[1mm] \texttt{MAX} y \texttt{MIN} \textbf{→} los extremos. \\[1mm] \texttt{CONTAR} \textbf{→} cuántos casos. \\[1mm] \textbf{Exigen:} una columna de un solo tipo y un rango completo.
\end{drawbox}

\begin{softbox}{milcMostaza}{milcGris}{Errores comunes y su corrección}
(1) \textbf{Seleccionar celda por celda}: la fórmula deja de crecer con la tabla. (2) \textbf{Promediar una columna con texto adentro}: devuelve un número, y es falso. (3) \textbf{Quedarse solo con el promedio}: esconde a los extremos, que suelen ser lo importante. (4) \textbf{Interpretar en abstracto}: «el promedio es 4,2» no dice nada sin decir 4,2 de qué y para qué. (5) \textbf{No verificar}: un total que no cuadra con una suma a mano de tres filas ya avisa de algo.
\end{softbox}

\begin{infoband}{milcTurquesa}


\textbf{Lo que va al cuaderno (Actividad 2 · EXPLICA).} \textbf{Título:} «Actividad 2 --- Qué pregunta responde cada una». \textbf{Formato:} las cinco fórmulas con su pregunta y lo que exigen, más el caso donde el promedio esconde algo y el indicador que lo mostraría. \textbf{Extensión:} media página. \textbf{Sabes que terminaste cuando} tienen escrito un caso donde el promedio engaña.
\end{infoband}

\puente{Tienes las cinco. Toca aplicarlas a tu tabla y decir qué significa cada número en tu caso, no en abstracto.}

\milcestacion{milcMagenta}{Estación 3 de 3 · Hacer}

\begin{softbox}{milcMagenta}{milcGris}{Cómo vas a construir tu producto}
\textbf{Actividad 3 · APLICA --- Cinco indicadores interpretados} (30 min · individual). (1) Sobre tu tabla, escribe las cinco fórmulas con su sintaxis completa. (2) Comprueba que cada rango tome la columna entera. (3) Verifica al menos un resultado a mano, con tres filas. (4) Escribe la interpretación de cada indicador en una frase de tu contexto. (5) Señala cuál de los cinco sirve para la decisión que tenías en mente. (6) Pásale la hoja a un compañero y pídele que diga qué significa uno de tus números.
\end{softbox}

\milcpaso{1}{milcMagenta}{\textbf{Tu entrega tiene tres piezas.} Las cinco fórmulas con su sintaxis. Las cinco interpretaciones en una frase. Y la señal en el indicador que de verdad sirve.}
\milcpaso{2}{milcMagenta}{\textbf{Interpretar es hablar de tu tema.} «El máximo es 11» no interpreta nada. «El máximo es 11 horas, y es un compañero que llega sin dormir» sí, porque cualquiera entiende qué hacer con eso.}
\milcpaso{3}{milcMagenta}{\textbf{Verificar cuesta un minuto.} Suma tres filas a mano y compáralas con la fórmula. Si no coinciden, el problema está en el rango o en el tipo de la columna, y siempre es mejor descubrirlo ahora.}
\milcpaso{4}{milcMagenta}{\textbf{Uno de los cinco manda.} Los cinco se calculan, pero solo uno responde tu pregunta original. Señalarlo obliga a recordar para qué recogiste los datos.}

\begin{drawbox}{milcMagenta}{Antes de entregar}
\checkbox Las cinco fórmulas con su sintaxis completa. \\[1mm] \checkbox Todos los rangos toman la columna entera. \\[1mm] \checkbox Al menos un resultado verificado a mano. \\[1mm] \checkbox Cinco interpretaciones de una frase, en tu contexto. \\[1mm] \checkbox Señalado el indicador que sirve para tu decisión. \\[1mm] \checkbox Un compañero explicó uno de tus números sin ayuda.
\end{drawbox}

\begin{softbox}{milcMagenta}{milcGris}{Plantilla de trabajo}
\textbf{Indicador.} \textit{Cuál de los cinco.} \\[1mm] \textbf{Fórmula.} \textit{La sintaxis, con su rango.} \\[1mm] \textbf{Resultado.} \textit{El número.} \\[1mm] \textbf{Qué significa.} \textit{Una frase de mi contexto.} \\[1mm] \textbf{Sirve para.} \textit{La decisión que tenía en mente.}
\end{softbox}

\begin{infoband}{milcMagenta}


\textbf{Lo que va al cuaderno (Actividad 3 · APLICA).} \textbf{Título:} «Actividad 3 --- Cinco indicadores interpretados». \textbf{Formato:} las cinco fórmulas con su resultado y su interpretación de una frase, más la señal en el que sirve para tu decisión. \textbf{Extensión:} media página. \textbf{Sabes que terminaste cuando} un compañero explicó uno de tus números sin que le ayudaras.
\end{infoband}

\puente{Lo que entregas hoy: cinco indicadores con su interpretación. La prueba de calidad: cada frase habla de tu tema, no de estadística en general.}

\milcestacion{milcMostaza}{Producto de la sesión}
\textbf{Cinco indicadores calculados con fórmula e interpretados en contexto}.
\begin{softbox}{milcMostaza}{milcGris}{Criterios de calidad}
(1) Las cinco fórmulas están escritas con su sintaxis y su rango completo. (2) Al menos un resultado está verificado a mano. (3) Las cinco interpretaciones hablan del tema propio, no de estadística en general. (4) Está señalado cuál indicador sirve para la decisión original. (5) Un compañero pudo explicar uno de los números sin ayuda.
\end{softbox}

\puente{Antes de cerrar, mira lo que hiciste desde cinco lados. Un número sin interpretación es un número; con ella empieza a servir para algo.}

\milcestacion{milcVino}{Evaluación liberadora · cinco dimensiones}

\begin{softbox}{milcVino}{milcGris}{Sentido de la evaluación}
Evaluar no es solo revisar una nota. Es reconocer qué aprendí, qué cambió en mí, qué emociones manejé, cómo me relaciono con la ciudad y la comunidad, y qué le dejaré a quien viene detrás.
\end{softbox}

\textbf{Mi reflexión liberadora en cinco dimensiones.} Completa cada frase iniciada:

\small
\begin{tabularx}{\textwidth}{>{\bfseries\RaggedRight\arraybackslash}p{3.4cm}Y>{\RaggedRight\arraybackslash}p{4.6cm}}
\rowcolor{milcVino}\color{white}Dimensión & \color{white}\bfseries Mi respuesta (frame starter) & \color{white}\bfseries Algo que descubrí\\
1. Personal & Lo que más cambió en mí fue \linead{6.5cm} & \linead{4.2cm}\\
2. Emocional & La emoción más fuerte fue \linead{2.5cm} y la regulé \linead{3cm} & \linead{4.2cm}\\
3. Ciudadana & En democracia esto importa porque \linead{6cm} & \linead{4.2cm}\\
4. Local (Cartago) & En mi barrio sería útil para \linead{6.5cm} & \linead{4.2cm}\\
5. Intergeneracional & A mi abuela/abuelo le diría \linead{6.5cm} & \linead{4.2cm}\\
\end{tabularx}
\normalsize

\begin{infoband}{milcVino}
\textbf{Para la próxima sustentación.} Vuelve a esta tabla en seis meses. Verás que las respuestas cambian — no porque lo aprendido cambie, sino porque tú cambias. La evaluación liberadora es una conversación contigo a través del tiempo.
\end{infoband}


\puente{Para terminar, tres citas verificadas sobre calcular con propósito: Dussel sobre la técnica que echa mano de la ciencia, Marco Aurelio sobre lo que basta, Floridi sobre en qué proyecto estamos trabajando.}

\milcestacion{milcGradeDark}{Triángulo de pensamiento · cierre filosófico}

\begin{softbox}{milcVino}{milcGris}{Dussel · lente del nosotros}
\textit{«Es un error pensar que la ciencia pura… aplica ella misma alguna de sus conclusiones, apareciendo la tecnología como su concreta creación… Muy por el contrario, es el discurso técnico artesanal o tecnológico… echa mano… de conclusiones o teorías científicas.»}

{\footnotesize\itshape\color{milcNegro!70}--- Enrique Dussel · Filosofía de la liberación (1977), §5.5.2}

La tendera no aprendió a sumar porque alguien le enseñara estadística. Sumó porque hacía falta llevar la cuenta. Las fórmulas de hoy son eso mismo: una necesidad práctica que después alguien formalizó.

\textbf{¿Qué cálculo hago por necesidad, sin haberlo aprendido en una clase?}
\end{softbox}
\linead{\linewidth}\\[1mm]
\linead{\linewidth}

\begin{softbox}{milcMostaza}{milcGris}{Estoicismo · Marco Aurelio · Meditaciones IX, 6 (c. 175 d.C.) · cuidado interior}
\textit{«Tu opinión presente fundada en el entendimiento, tu conducta presente dirigida al bien común y tu disposición presente de contento con todo lo que sucede: eso basta. (trad. propia)»}

{\footnotesize\itshape\color{milcNegro!70}--- Marco Aurelio · Meditaciones IX, 6 (c. 175 d.C.)}

Cinco indicadores bastan para casi cualquier tabla de este tamaño. La tentación es calcular quince porque el programa los ofrece, y quince números sin interpretar informan menos que cinco bien leídos.

\textbf{¿Cuántos de los números que calculé sé explicar sin volver a mirarlos?}
\end{softbox}
\linead{\linewidth}\\[1mm]
\linead{\linewidth}

\begin{softbox}{milcTurquesa}{milcGris}{Floridi · lente de la infoesfera}
\textit{«Una de las preguntas políticas apremiantes que enfrentamos en las sociedades de la información avanzadas es: ¿en qué clase de proyecto humano estamos trabajando? (trad. propia)»}

{\footnotesize\itshape\color{milcNegro!70}--- Luciano Floridi · Commentary on the Onlife Manifesto (2015), § 3.1}

La versión pequeña de esa pregunta cabe en tu hoja: para qué estás calculando. Un promedio sin pregunta detrás es aritmética; con pregunta detrás es el comienzo de una decisión.

\textbf{¿Para qué recogí estos datos, y sigo calculando lo que responde a eso?}
\end{softbox}
\linead{\linewidth}\\[1mm]
\linead{\linewidth}

\begin{infoband}{milcNegro}
\textbf{Nota para el docente.} Las tres son citas textuales y llevan su referencia debajo. Puedes remitir a la obra en clase.
\end{infoband}

\milcestacion{milcVino}{Mi autoevaluación · marca una casilla por fila}

{\small
\setlength{\extrarowheight}{2.2pt}
\begin{tabularx}{\textwidth}{>{\RaggedRight\arraybackslash\bfseries}p{3.0cm}YYY}
\rowcolor{milcVino}\color{white}\bfseries Criterio & \color{white}\bfseries Lo logré & \color{white}\bfseries En proceso & \color{white}\bfseries Necesito apoyo\\[.6mm]
Comprensión del tema & \milcopcion{Explico las ideas con mis palabras.} & \milcopcion{Reconozco las ideas, falta precisión.} & \milcopcion{Necesito repasar lo básico.}\\[.8mm]
\rowcolor{milcGris}Pensamiento computacional & \milcopcion{Usé los pilares al armar mi producto.} & \milcopcion{Usé algunos pilares.} & \milcopcion{Necesito releer la estación 2.}\\[.8mm]
Aplicación y producto & \milcopcion{Mi producto cumple los criterios.} & \milcopcion{Está armado, con ajustes pendientes.} & \milcopcion{Aún está en construcción.}\\[.8mm]
\rowcolor{milcGris}Reflexión 5 dimensiones & \milcopcion{Respondí cada una con honestidad.} & \milcopcion{Respondí algunas con detalle.} & \milcopcion{Mis respuestas son muy generales.}\\[.8mm]
Triángulo de pensamiento & \milcopcion{Conecté las 3 voces con mi trabajo.} & \milcopcion{Conecté al menos una.} & \milcopcion{No las relaciono con mi proceso.}\\[.8mm]
\end{tabularx}}

\vspace{3mm}
\textbf{Hoy entendí que:}\\[1mm]
\linead{\linewidth}\\[1mm]
\linead{\linewidth}

\textbf{Mi compromiso para la próxima sesión es:}\\[1mm]
\linead{\linewidth}\\[1mm]
\linead{\linewidth}

\begin{softbox}{milcMostaza}{milcGris}{Cita de despedida · Marco Aurelio}
\textit{``El obstáculo en el camino se vuelve el camino. Lo que estorba al camino se vuelve camino.''}\\[1mm]
Cada error de hoy, bien observado, se convierte en pieza del camino que pisarás mañana.
\end{softbox}

\small
\begin{tabularx}{\textwidth}{>{\bfseries}p{3.6cm}Y}
\rowcolor{milcGradeDark!12}Nombre completo: & \linead{10cm}\\
Fecha: & \linead{4cm} \quad Curso/grupo: \linead{4cm}\\
Mi firma: & \linead{9cm}\\
\end{tabularx}
\normalsize

\begin{infoband}{milcGradeDark}
\textbf{Para la próxima sesión.} Llega con los cinco indicadores interpretados. Pregúntale a alguien que lleve cuentas a mano, en una tienda o en la casa, qué revisa cuando cuadra. En la sesión 4 aprendes a preguntarle a la tabla solo por una parte.
\end{infoband}

\end{document}
